\chapter{Sintonización clásica del PID}
\label{ch:sintonizacion-clasica}

\section{Introducción del capítulo}
La historia del controlador PID está estrechamente ligada a métodos de ajuste prácticos que
buscaban parametrizar el regulador con rapidez y sin necesidad de resolver problemas de
optimización complejos. En muchos entornos industriales, estas reglas siguen siendo el punto
de partida porque son fáciles de aplicar y permiten obtener respuestas razonables en procesos
moderados. Sin embargo, su validez depende de hipótesis que no siempre se cumplen en
sistemas con baja inercia, dinámica hidráulica o perturbaciones estocásticas.

\section{Objetivo del capítulo}
Este capítulo describe los métodos clásicos de sintonización PID y discute sus fundamentos,
supuestos y limitaciones. La intención no es descartarlos por completo, sino mostrar en qué
condiciones resultan útiles, cuándo se vuelven insuficientes y por qué sirven como línea base
adecuada para comparar técnicas de ajuste basadas en metaheurísticas.

\section{Métodos empíricos}
Los métodos empíricos se construyeron a partir de observaciones experimentales y reglas
heurísticas aplicadas a familias de procesos. Su principal fortaleza es la rapidez: con unas
pocas pruebas de lazo abierto o lazo cerrado es posible proponer valores iniciales para
$K_p$, $K_i$ y $K_d$. Entre sus debilidades destacan la dependencia del operador, la dificultad
para reproducir resultados y la escasa adaptabilidad cuando la planta presenta no linealidades
o cambia de régimen de operación.

\section{Método de Ziegler--Nichols}
El método de Ziegler--Nichols, en su versión de lazo cerrado, consiste en llevar el sistema al
límite de estabilidad hasta identificar la ganancia última $K_u$ y el período último $T_u$. A
partir de esos dos valores, una tabla fija propone las ganancias del controlador. La lógica del
método es simple y poderosa: si el proceso puede describirse por su oscilación crítica,
entonces es posible obtener una sintonización rápida sin resolver un modelo detallado. No
obstante, esa simplicidad suele traducirse en ajustes agresivos cuando la planta no responde
como un proceso lineal bien amortiguado.

\section{Otros métodos clásicos}
Además de Ziegler--Nichols, la literatura clásica incluye enfoques como Cohen--Coon, reglas
de reacción del proceso y procedimientos iterativos de prueba y error. Cohen--Coon mejora el
ajuste en ciertos sistemas con retardo significativo, mientras que la prueba y error permite al
operador incorporar experiencia práctica sobre la planta. A pesar de esas variantes, todos
comparten una limitación de fondo: su diseño se apoya en aproximaciones locales y no en una
exploración global del espacio de soluciones.

\section{Limitaciones de los métodos clásicos}
En sistemas como la microred de Limbani, donde intervienen una turbina Michell--Banki,
efectos de fase no mínima y perturbaciones rápidas por variación de irradiancia, las reglas
clásicas pueden conducir a respuestas subamortiguadas y oscilatorias. El controlador logra
reaccionar, pero lo hace con un costo elevado en tiempo de asentamiento, esfuerzo del
servomecanismo y calidad de la frecuencia. Estas limitaciones justifican el tránsito hacia
métodos de optimización que no dependan exclusivamente de una fórmula cerrada, sino de la
evaluación repetida del desempeño real del sistema.

\section{Cierre del capítulo}
La sintonización clásica conserva valor como referencia práctica y como punto de comparación.
Sin embargo, cuando el problema requiere considerar múltiples restricciones y objetivos
dinámicos al mismo tiempo, resulta más natural reformular el ajuste como un problema de
optimización. Esa reformulación se desarrolla en el siguiente capítulo.
